\section{Análisis por Componente: Rejillas, Damper, Instrumentación y Accesorios}

\subsection{Rejillas de Exfiltración}

Tres rejillas de 353 mm $\times$ 336 mm (área facial unitaria 0.1187 m$^2$) descargan el caudal a velocidad facial de 3 m/s con $\Delta P$ de 11 Pa en el sitio (cierre de orificio, $C_d$ = 0.60). El tipo indicado para máxima área libre es eggcrate (panal) 1/2$\times$1/2$\times$1/2 in, con área libre de aproximadamente 90 \% según catálogo \cite{airfoil_rcfcr5}; la malla anti-insectos inox 316 de 18$\times$18 aporta aproximadamente 48-51 \% de área abierta \cite{imm_malla}. El área neta combinada ($\approx$45 \% del área facial, dato típico) es la base del cálculo de la memoria; el comportamiento de las rejillas de transferencia como elementos de alivio de presión está documentado en la literatura técnica \cite{bsc_ba0006}. La Tabla \ref{tab:inf002_rejillas} presenta la comparativa de candidatas.

\begin{table}[htbp]
	\centering
	\caption{Comparativa de rejillas candidatas (consulta 23/07/2026)}
	\label{tab:inf002_rejillas}
	\setcounter{itemcount}{0}
	\footnotesize
	\begin{tabularx}{\textwidth}{@{}NXlXXXc@{}}
		\toprule
		\multicolumn{1}{c}{\textbf{It.}} & \textbf{Producto} & \textbf{Tipo} & \textbf{Área Libre} & \textbf{Materiales} & \textbf{Pros / Contras} & \textbf{Fuente} \\
		\midrule
		& Titus 50F / 50R & Eggcrate retorno 1/2 in & La mayor área libre de su categoría ($\sim$90 \%, dato típico) & Aluminio; versión íntegramente inox disponible & Única con ejecución inox de fábrica; importación & \cite{titus_50f} \\
		& Krueger EGC5 & Eggcrate 1/2 in & Maximiza el área libre & Aluminio & Amplia gama de tamaños; no inox & \cite{krueger_egc5} \\
		& Laminaire (Colombia) & Rejillas y difusores a pedido & Según diseño & Aluminio; fabricación local a medida & Permite 353 mm $\times$ 336 mm exacto y marco para malla; confirmar ejecución inox & \cite{laminaire} \\
		& ProAire S.A.S (Bogotá) & Rejillas con corte láser a medida & Según diseño & Inox 316 fabricable bajo pedido & Vía más robusta para inox 316L local & \cite{proaire} \\
		\bottomrule
	\end{tabularx}
\end{table}

La decisión es: rejilla eggcrate inox 316L (o aluminio anodizado como alternativa interior) con malla anti-insectos inox 316 18$\times$18 montada en marco desmontable para limpieza; aislamiento del par galvánico malla-inox/marco-aluminio si se mezclan materiales (dato típico de ingeniería).

\subsection{Damper de Alivio Barométrico}

El damper barométrico es un backdraft damper con presión de inicio de apertura ajustable por contrapeso; es la solución pasiva estándar para controlar presurización positiva sin controles activos \cite{greenheck_backdraft}. El set-point de +25 Pa $\approx$ 0.10 in c.a. cae dentro del rango de ajuste de los candidatos, como se aprecia en la Tabla \ref{tab:inf002_dampers}.

\begin{table}[htbp]
	\centering
	\caption{Comparativa de dampers de alivio (consulta 23/07/2026)}
	\label{tab:inf002_dampers}
	\setcounter{itemcount}{0}
	\footnotesize
	\begin{tabularx}{\textwidth}{@{}NXXXXc@{}}
		\toprule
		\multicolumn{1}{c}{\textbf{It.}} & \textbf{Fabricante / Modelo} & \textbf{Rango de Ajuste Start-Open} & \textbf{Materiales} & \textbf{Pros / Contras} & \textbf{Fuente} \\
		\midrule
		& Greenheck BR-40/41/42 y SEBR-10 (severe environment) & Seleccionable 0.05--0.30 in c.a. (12--75 Pa); cubre +25 Pa & SEBR: construcción ambiente severo & Línea severe-environment pertinente para hipoclorito; canal Greenheck en Bogotá & \cite{greenheck_sebr10} \\
		& Ruskin CBD6 / BD6 & Control de presión estática hasta 0.25 in c.a. (62 Pa) & Aluminio extruido 6063-T5 resistente a corrosión & Todo aluminio, idóneo para cloro; el set-point queda cerca del límite inferior del rango & \cite{ruskin_bd6} \\
		& Nailor 1390CB & Contrapeso externo 360$^{\circ}$; alivio a diferenciales extremadamente bajos & Opción aluminio extruido 6063-T5; sellos neopreno & Rodamientos de bolas: máxima sensibilidad a bajos diferenciales; importación & \cite{nailor_1390cb} \\
		& Halton BRD & Apertura mínima 30--200 Pa; ajustable hasta 300 Pa & Construcción marina & Sobrespecificado en presión; el mínimo de 30 Pa excede el set-point de 25 Pa & \cite{halton_brd} \\
		\bottomrule
	\end{tabularx}
\end{table}

La decisión es: damper barométrico con contrapeso ajustable en el rango 12-75 Pa, construcción aluminio extruido 6063-T5 o línea severe-environment (Greenheck SEBR-10 como primera opción por canal local; Ruskin CBD6 y Nailor 1390CB-EAF como alternativas), con ensayos AMCA 500-D y calibración final en comisionamiento. Se descartan marcos galvanizados estándar. El damper se dimensiona para el caudal de alivio de diseño con velocidad facial $\leq$ 2.5 m/s (dato típico, baja pérdida y sin ruido).

\subsection{Instrumentación de Presión Diferencial}

La Tabla \ref{tab:inf002_instrumentos} presenta la comparativa de instrumentos candidatos para el lazo de presión diferencial.

\begin{table}[htbp]
	\centering
	\caption{Comparativa de instrumentos candidatos (consulta 23/07/2026)}
	\label{tab:inf002_instrumentos}
	\setcounter{itemcount}{0}
	\footnotesize
	\begin{tabularx}{\textwidth}{@{}NXXXXc@{}}
		\toprule
		\multicolumn{1}{c}{\textbf{It.}} & \textbf{Instrumento} & \textbf{Rango / Salida} & \textbf{Precisión / Protección} & \textbf{Pros / Contras} & \textbf{Fuente} \\
		\midrule
		& Dwyer Magnehelic 2000-00 (indicador local) & 0--0.25 in c.a. (0--62 Pa) & $\pm$2 \% FS; caja aluminio con ensayo niebla salina 168 h; opción bisel inox & Estándar de facto; set-point +25 Pa al 40 \% de escala; distribuidores en Colombia & \cite{skilltech_magnehelic} \\
		& Dwyer MS-121(-LCD) Magnesense (transmisor) & 0.1/0.25/0.5 in c.a. seleccionables (25/62.5/125 Pa); 4-20 mA 2 hilos & $\pm$1 \% FS; NEMA 4X (IP66); constante de tiempo ajustable 0.5--15 s & Primera opción: protección IP66 crítica en ambiente clorado; amortigua ráfagas de viento & \cite{dwyer_ms} \\
		& Setra 264 (transmisor alternativo) & 0--0.25 in c.a. (0--62 Pa); 4-20 mA & $\pm$0.25/$\pm$0.4/$\pm$1 \% FS; sobrepresión 10 psi & Elemento capacitivo inox, alta estabilidad; suministro por importación (por confirmar stock) & \cite{setra_264} \\
		& Siemens QBM2130-1U (transmisor alternativo) & 0--100 Pa; 4-20 mA & IP42 & Viable si el BMS es Siemens; IP42 exige gabinete adicional & \cite{siemens_qbm2130} \\
		& Dwyer Photohelic 3000MR-00AV (indicador + alarmas) & 0--0.25 in c.a.; 2 relés SPDT & Repetibilidad switch 1 \% & Genera alarma alta/baja sin PLC si el cliente no tiene BMS & \cite{photohelic_nci} \\
		\bottomrule
	\end{tabularx}
\end{table}

La arquitectura del lazo (detalle en HD-INST-001) es: transmisor MS-121 entre el interior y una referencia exterior protegida (caja estática de pared orientada a sotavento), alarma baja +12.5 Pa, alarma alta +40 Pa, retardo 30-60 s; Magnehelic 2000-00 de lectura local; damper barométrico como elemento final pasivo. El lazo activo con variador de frecuencia solo se justifica si hay aperturas frecuentes de puerta o exigencia de registro BMS (dato típico de ingeniería \cite{aivc_7469}). La disponibilidad Dwyer en Colombia comprende Vía Industrial (Bogotá/Cali), RS Importaciones y Suministros, y SICO Global \cite{viaindustrial_dwyer}.

\subsection{Accesorios de Montaje}

La ferretería y los soportes se especifican en inox AISI 316 / clase A4 como mínimo: el molibdeno mejora la resistencia a picadura y corrosión por rendija frente a cloruros \cite{marsh_fasteners}, con aislamiento dieléctrico frente a estructura galvanizada. Los soportes estándar zincados de los instrumentos se sustituyen por fabricación local en inox 316; los cuerpos de aluminio fundido de los instrumentos no se montan a la intemperie.

Los sellantes se especifican en silicona RTV neutra como primera opción, compatible con hipoclorito \cite{rtv_silicone}; se evita la silicona acetoxi (libera ácido acético corrosivo) y el poliuretano en exposición directa a cloro gaseoso \cite{apache_pu}.

La conexión flexible del ventilador se especifica en hipalón (CSM), material recomendado para exterior por resistencia química, ozono y envejecimiento \cite{hardcast_hypalon}, con bandas de amarre inox 316 en lugar del galvanizado estándar.
